\chapternotnumbered{Introduction} \label{ch:Introduction}

Modern numerical linear algebra is increasingly shaped by two practical facts. The first is that many linear operators arising in scientific computing, statistics, and data analysis are represented by matrices so large that classical dense algorithms become too expensive to apply without modification. The second is that modern hardware no longer offers a single uniform arithmetic environment: computation is distributed across several precisions, often with large differences in throughput, memory footprint, and numerical reliability.

Randomized numerical linear algebra addresses the first point by replacing expensive deterministic steps with carefully designed randomized approximations. Typical examples include sketching, random projections, and structured transforms that preserve enough geometric information about a relevant subspace to make the resulting algorithms both faster and more scalable. These ideas are especially attractive for overdetermined least-squares problems, where one seeks a good approximation to a vector in the range of a linear map while avoiding the full cost of the original problem.

Mixed-precision arithmetic addresses the second point. Instead of performing all operations in the same format, one attempts to use lower precision where it is computationally advantageous and higher precision only where it is numerically necessary. The resulting tradeoff is subtle: lower precision may reduce runtime and memory traffic, but it can also enlarge perturbations, worsen the effective conditioning of intermediate algebraic problems, or even destabilize an otherwise sound randomized method.

This thesis is motivated by the interaction of these two ideas. Randomized algorithms are attractive because they reorganize computations into operations that can be implemented efficiently on modern hardware. Mixed precision is attractive because it offers a direct route to further computational gains. However, the combination is not automatically safe: sketching modifies the geometry of the problem by replacing one operator with a compressed surrogate, while low precision modifies the arithmetic through additional rounding perturbations. Understanding how these two approximations interact is therefore one of the central goals of this work.

The main mathematical focus is on randomized methods for linear least-squares problems, especially sketching with Gaussian and structured random transforms such as the subsampled randomized Hadamard transform. The thesis studies the quality of the resulting approximations, the behavior of residual and solution errors, and the role of conditioning in the reduced systems. These questions are then revisited in the presence of mixed-precision arithmetic, where different parts of the algorithm may be computed in different formats and therefore subjected to different perturbation models.

The climate-model material appears later in the thesis as an application rather than as the conceptual center. In particular, the goal is not to write a climate-science thesis, but to test whether randomized and mixed-precision methods developed on the mathematical side remain useful in a realistic downstream workflow. Even in that setting, useful outcomes are not limited to raw speedup; reduced memory use, acceptable lower-precision behavior, or a clearer understanding of failure regimes are also valuable.

The structure of the thesis is as follows. The preliminary chapters fix notation and collect the mathematical background needed for sketching methods and floating-point analysis. The next chapters discuss mixed-precision arithmetic and randomized least squares in more detail, followed by numerical experiments designed to measure accuracy, conditioning, and computational tradeoffs. The final substantive chapter considers the climate-model application and places the earlier analysis into a more realistic computational context.

\begin{remark}[Current Stage]
    At the current stage of writing, the mathematical background and theoretical organization are still being developed. For that reason, some later chapters are presently lightweight scaffolds whose role is to establish the eventual structure of the manuscript.
\end{remark}
